\section{Threats to Validity}
\label{sec:threats}

\paragraph{Simulator abstraction.} All results in this manuscript are
produced by a discrete-event simulator, not a real serving engine
(\S\ref{sec:execution-model}). We do not claim that simulated absolute
latency or throughput values match real-hardware values; only relative
ranking and reversal direction are claimed to transfer, and only for the
frozen validation cases in \S\ref{sec:real-system}, once that section's
result exists.

\paragraph{Fidelity of reimplemented policies.} \texttt{FAITHFUL\_EXTERNAL}
policies are validated reimplementations, not the original authors' code;
residual implementation gaps relative to the original published systems
are possible despite unit-test validation (\S\ref{sec:fidelity}).

\paragraph{\texttt{STYLE\_APPROXIMATION} baselines.} The two style-level
approximation policies are explicitly excluded from the PRIMARY headline
panel and used only for a robustness check; they must not be read as
official or validated reimplementations of their namesake systems
(\S\ref{sec:fidelity}).

\paragraph{Three primary workload families.} The Pilot-V2 corpus draws from
three primary sources (\S\ref{sec:sources}); a fourth candidate source
(TraceLab) is retained only as out-of-distribution context because its
independence from a prior public window sweep cannot be fully verified.
Any portability finding in \S\ref{sec:results} is scoped to these three
sources plus whatever out-of-distribution evidence TraceLab can
responsibly contribute, not to LLM-serving workloads in general.

\paragraph{Dataset/provider dependence.} As established in
\S\ref{sec:sources}, BurstGPT and this project's Azure traces are not
established to come from independent underlying providers -- both trace
back to Microsoft Azure infrastructure, collected through different
mechanisms. Any cross-source finding involving BurstGPT and Azure should be
read with this in mind; Bailian/Qwen remains a genuinely independent
platform.

\paragraph{Prompt-token semantic differences.} Prompt-token fields compose
differently across sources (single-request figure for BurstGPT versus
potentially multi-turn-cumulative context for Azure and TraceLab;
\S\ref{sec:descriptors}). Any prompt-length-based claim in this manuscript
carries this caveat and should not be read as evidence that the sources'
underlying conversations differ in length, only that their observed,
request-level prompt-token figures do.

\paragraph{Unresolved historical byte-level overlap.} An exact,
byte-level reconstruction of a prior related project's specific 60-window
public-trace replay set could not be independently re-derived from this
project's available checkout; that project's own recorded parameters and
final verdict are taken on the strength of direct manuscript inspection,
not independently re-derived from local artifacts. This project's own
window construction is independently verified as non-identical
(\S\ref{sec:windows}, \S\ref{sec:evidence-boundaries}), but the historical
byte-level reconstruction gap itself remains an open item.

\paragraph{Homogeneous hardware / execution assumptions.} The simulator
executes every cell under a fixed, single-GPU-colocated execution model for
the PRIMARY panel; two disaggregation-requiring policies are confined to a
secondary stratum and never pooled with the PRIMARY leaderboard
(\S\ref{sec:panel}).

\paragraph{Completion-conditioned undefined metrics.} Several metrics are
undefined, not zero, when a cell records zero completions
(\S\ref{sec:metrics}); this is a real methodological choice with a
real cost -- an undefined-metric policy is excluded from that specific
condition's ranking-derived statistics, which lowers effective sample size
for those conditions and must be read alongside every ranking-derived
table's reported exclusion count.

\paragraph{Phase-11 zero-completion prevalence.} 301 of 720 FIFO-only
calibration-region records recorded zero completions
(\S\ref{sec:calibration}). This is disclosed as a calibration
characteristic, not a headline finding; it does, however, mean that a
non-trivial share of (window, region) combinations will contribute no
completion-conditioned-metric observation to \S\ref{sec:results}'s
downstream analyses for any policy at that specific region, which is a
structural constraint on statistical power in the highest-pressure regions
in particular.

\paragraph{SLO-definition sensitivity is disclosed, not claimed.} An
SLO-synthesis sensitivity check was planned but no alternative
SLO-synthesis rule was ever actually preregistered before or during
Pilot-V2 execution (\S\ref{sec:slo-sensitivity-gap}); no alternative rule
was invented after the fact, no such experiment has been run, and this
robustness component is not claimed as evidence anywhere in this
manuscript. The mitigation is disclosure and process, not substitution: any
future SLO-sensitivity study will be run as an explicitly labeled
post-campaign extension with its own frozen protocol, never retrofitted
into this study's preregistered evidence.

\paragraph{Scope of real-system validation.} \S\ref{sec:real-system}'s
validation is limited by design to a small, frozen set of case-selected
conditions (largest-effect-size reversal plus one stable-ordering control),
not a full re-execution of the Pilot-V2 matrix on real hardware; agreement
or disagreement on those specific cases should not be over-generalized to
every cell in \S\ref{sec:results}.
